% GENERATED FILE. Edit canonical lessons, then run npm run generate:books.
% canonical-source-hash: fnv1a64:fd2d06f7a8a418dd
% canonical-lessons: ZH-C477-xinzang, ZH-C477-gutou, ZH-C477-danao, ZH-C477-changwei, ZH-C477-weikou

\chapter{Inside the Body}
\label{ch:zh-inside-the-body}
\hlchaptermodality{477}

\begin{chapteropening}
\textbf{By the end of this chapter:} \emph{I can talk about the heart, bones and the brain, my stomach and my appetite.}

Complete the last lesson of chapter 477: I can talk about the heart, bones and the brain, my stomach and my appetite.
\end{chapteropening}

\section[xīnzàng]{xīnzàng — the heart as an organ}

\label{lesson:ZH-C477-xinzang}



\begin{quote}
Before the new one: say the Mandarin for to shake one's head, then the Mandarin for bodily pain.
\end{quote}

\subsection*{You'll want to know: xīnzàng}
\emph{xīnzàng} — \textquotedblleft{}the heart as an organ\textquotedblright{}.

\begin{quote}
\emph{wǒ yéye de xīnzàng bú tài hǎo} — My grandfather has a weak heart.
\end{quote}

Its characters come later; for now it is a word for the ear.

\begin{cousinweb}
\emph{xīnzàng} is built from \zh{心} 'heart' and \emph{zàng} 'an inner organ'. The first is inherited from Middle Chinese \emph{sim} (Guangyun, in Baxter's transcription), glossed as the organ of fire; no dictionary we use records an older form of the second under this spelling, so its history is not settled.
\end{cousinweb}

\begin{grammarlens}[title={What you've built}]
One more word for saying how it was.
\end{grammarlens}

\subsection*{Your turn}
\begin{itemize}
\raggedright
  \item \emph{Say it:} \emph{xīnzàng}
  \item \emph{Say it:} \emph{xīnzàng}, once more
  \item \emph{Say it:} \emph{wǒ yéye de xīnzàng bú tài hǎo} — My grandfather has a weak heart.
  \item \emph{Recall:} say the Mandarin for to shake one's head, then the Mandarin for bodily pain, then say \emph{xīnzàng} again
\end{itemize}

\begin{tcolorbox}[breakable,colback=teal!4,colframe=teal!35!black,title={Before you move on}]
What does xīnzàng mean? (\textquotedblleft{}The heart as an organ\textquotedblright{}.) Say it once more.
\end{tcolorbox}
\section[gǔtou]{gǔtou — bone}

\label{lesson:ZH-C477-gutou}



\begin{quote}
Before the new one: say the Mandarin for bodily pain, then the Mandarin for the heart as an organ.
\end{quote}

\subsection*{You'll want to know: gǔtou}
\emph{gǔtou} — \textquotedblleft{}bone\textquotedblright{}.

\begin{quote}
\emph{gǒu zuì xǐhuan chī gǔtou} — Dogs love chewing bones more than anything.
\end{quote}

Its characters come later; for now it is a word for the ear.

\begin{cousinweb}
\emph{gǔtou} is built from \emph{gǔ} 'bone' and \zh{头}, a noun ending here, both inherited from Middle Chinese, \emph{kwot} and \emph{duw} (Guangyun, in Baxter's transcription), which explains the first, after an older book, as the core of the flesh; the second is said light.
\end{cousinweb}

\begin{grammarlens}[title={What you've built}]
One more word for saying how it was.
\end{grammarlens}

\subsection*{Your turn}
\begin{itemize}
\raggedright
  \item \emph{Say it:} \emph{gǔtou}
  \item \emph{Say it:} \emph{gǔtou}, once more
  \item \emph{Say it:} \emph{gǒu zuì xǐhuan chī gǔtou} — Dogs love chewing bones more than anything.
  \item \emph{Recall:} say the Mandarin for bodily pain, then the Mandarin for the heart as an organ, then say \emph{gǔtou} again
\end{itemize}

\begin{tcolorbox}[breakable,colback=teal!4,colframe=teal!35!black,title={Before you move on}]
What does gǔtou mean? (\textquotedblleft{}Bone\textquotedblright{}.) Say it once more.
\end{tcolorbox}
\section[dànǎo]{dànǎo — brain}

\label{lesson:ZH-C477-danao}



\begin{quote}
Before the new one: say the Mandarin for the heart as an organ, then the Mandarin for a bone.
\end{quote}

\subsection*{You'll want to know: dànǎo}
\emph{dànǎo} — \textquotedblleft{}brain\textquotedblright{}.

\begin{quote}
\emph{shuìjiào de shíhou dànǎo yě zài gōngzuò} — Even while you sleep, your brain is working.
\end{quote}

Its characters come later; for now it is a word for the ear.

\begin{cousinweb}
\emph{dànǎo} is built from \zh{大} 'big' and \emph{nǎo} 'the brain', both inherited from Middle Chinese, \emph{dajH} and \emph{nawX} (Guangyun, in Baxter's transcription), which glosses the second as the brain in the head.
\end{cousinweb}

\begin{grammarlens}[title={What you've built}]
One more word for saying how it was.
\end{grammarlens}

\subsection*{Your turn}
\begin{itemize}
\raggedright
  \item \emph{Say it:} \emph{dànǎo}
  \item \emph{Say it:} \emph{dànǎo}, once more
  \item \emph{Say it:} \emph{shuìjiào de shíhou dànǎo yě zài gōngzuò} — Even while you sleep, your brain is working.
  \item \emph{Recall:} say the Mandarin for the heart as an organ, then the Mandarin for a bone, then say \emph{dànǎo} again
\end{itemize}

\begin{tcolorbox}[breakable,colback=teal!4,colframe=teal!35!black,title={Before you move on}]
What does dànǎo mean? (\textquotedblleft{}Brain\textquotedblright{}.) Say it once more.
\end{tcolorbox}
\section[chángwèi]{chángwèi — stomach and gut; digestion}

\label{lesson:ZH-C477-changwei}



\begin{quote}
Before the new one: say the Mandarin for a bone, then the Mandarin for the brain.
\end{quote}

\subsection*{You'll want to know: chángwèi}
\emph{chángwèi} — \textquotedblleft{}stomach and gut; digestion\textquotedblright{}.

\begin{quote}
\emph{wǒ chángwèi bù shūfu le liǎng tiān} — My stomach has been upset for two days.
\end{quote}

Its characters come later; for now it is a word for the ear.

\begin{cousinweb}
\emph{chángwèi} is built from \emph{cháng} 'the gut' and \emph{wèi} 'the stomach', both inherited from Middle Chinese, \emph{drjang} and \emph{hjw+jH} (Guangyun, in Baxter's transcription), which glosses the first in this very pair and the second, after an older book, as the storehouse of grain.
\end{cousinweb}

\begin{grammarlens}[title={What you've built}]
One more word for saying how it was.
\end{grammarlens}

\subsection*{Your turn}
\begin{itemize}
\raggedright
  \item \emph{Say it:} \emph{chángwèi}
  \item \emph{Say it:} \emph{chángwèi}, once more
  \item \emph{Say it:} \emph{wǒ chángwèi bù shūfu le liǎng tiān} — My stomach has been upset for two days.
  \item \emph{Recall:} say the Mandarin for a bone, then the Mandarin for the brain, then say \emph{chángwèi} again
\end{itemize}

\begin{tcolorbox}[breakable,colback=teal!4,colframe=teal!35!black,title={Before you move on}]
What does chángwèi mean? (\textquotedblleft{}Stomach and gut; digestion\textquotedblright{}.) Say it once more.
\end{tcolorbox}
\section[wèikǒu]{wèikǒu — appetite}

\label{lesson:ZH-C477-weikou}



\begin{quote}
Before the new one: say the Mandarin for the brain, then the Mandarin for one's digestion.
\end{quote}

\subsection*{You'll want to know: wèikǒu}
\emph{wèikǒu} — \textquotedblleft{}appetite\textquotedblright{}.

\begin{quote}
\emph{wǒ bìng le sān tiān wèikǒu bù hǎo} — I've been ill for three days and have no appetite.
\end{quote}

Its characters come later; for now it is a word for the ear.

\begin{cousinweb}
\emph{wèikǒu} is built from \emph{wèi} 'the stomach' and \zh{口} 'mouth', both inherited from Middle Chinese, \emph{hjw+jH} and \emph{khuwX} (Guangyun, in Baxter's transcription); the dictionary explains the first, after an older book, as the storehouse of grain.
\end{cousinweb}

\begin{grammarlens}[title={What you've built}]
One more word for saying how it was.
\end{grammarlens}

\subsection*{Your turn}
\begin{itemize}
\raggedright
  \item \emph{Say it:} \emph{wèikǒu}
  \item \emph{Say it:} \emph{wèikǒu}, once more
  \item \emph{Say it:} \emph{wǒ bìng le sān tiān wèikǒu bù hǎo} — I've been ill for three days and have no appetite.
  \item \emph{Recall:} say the Mandarin for the brain, then the Mandarin for one's digestion, then say \emph{wèikǒu} again
\end{itemize}

\begin{tcolorbox}[breakable,colback=teal!4,colframe=teal!35!black,title={Before you move on}]
What does wèikǒu mean? (\textquotedblleft{}Appetite\textquotedblright{}.) Say it once more.
\end{tcolorbox}
